\section{Motivation}
\label{sec:motivation}

Modern key--value stores --- RocksDB, LevelDB, Cassandra, HBase --- are built
on the \emph{log-structured merge-tree} (LSM-tree)~\citep{oneil1996lsm}.
The LSM-tree is a write-optimised data structure: incoming writes accumulate in an in-memory
buffer and flush to disk as immutable Sorted String Tables (SSTables)
at the top level, while a background compaction process merges and promotes
files to successively larger levels.
The structure converts random writes into
sequential I/O at the cost of \emph{read amplification}: a single lookup
has to consult an SSTable at every level.

To avoid reading every SSTable during a lookup, every file carries a small in-memory
\emph{filter} in its metadata.
The filter answers a membership question
with one-sided error: false negatives are forbidden, and the false positive rate (FPR) is bounded by~$\varepsilon$.
If the filter indicates that a key is not present in an SSTable, the engine
skips the file with zero disk I/O.
The Bloom filter~\citep{bloom1970space}
fills this role for point lookups and is the default in production
LSM-trees~\citep{monkey2017,dong2021rocksdb}.
Answering a range query with a Bloom filter, however, requires probing every key in the interval individually --- $O(\mathcal{L})$ probes in the worst case, which is impractical at LSM scale~\citep{zhang2018surf}.
Range-shaped workloads --- short scans, prefix lookups over hierarchical keys (file paths, S2 cells,
URLs), and time-window queries --- are common in production LSM workloads~\citep{cao2020facebook}.

For range queries we need a \emph{range filter}: given a query interval
$[a, b]$ of length at most~$\mathcal{L}$, decide whether $[a, b] \cap S = \emptyset$
under the same one-sided error guarantee.
\citet{goswami2015approximate} prove a tight worst-case lower bound of
$\log_2(\mathcal{L}/\varepsilon)$ bits per key and give a matching
construction --- a locality-preserving hash followed by an
\emph{Exact Range Emptiness} (ERE) backend --- but their construction
is, in their successors' own words, ``mainly theoretic in nature and
thus very complicated to implement''~\citep{grafite2024}.

The bound is worst-case, however, and \citet{vaidya2022snarf} make
the escape clause explicit: it holds only \emph{without further assumptions on the data or workload},
so a filter tuned to a specific data or workload can do better.
Existing practical filters respond to this in two ways:
either go below the Goswami bound by assuming a specific data distribution,
losing the worst-case FPR guarantee (SuRF, SNARF),
or stay at the bound and guarantee the FPR under any workload
(Rosetta, Grafite, Memento)~\citep{luo2020rosetta,grafite2024,eslami2024memento}.
Many production LSM keys, however, have exploitable structure: they pack densely across the used range, spread across the universe, or concentrate in dense regions separated by large gaps~\citep{cao2020facebook}.
None of the filters above uses this.

This thesis decomposes, simplifies, and accelerates the Goswami-style
\emph{Exact Range Emptiness} (ERE) backend.
We implemented Goswami's internal Weak Prefix Search
(WPS) index~\citep{belazzougui2010fast} in full --- the part Grafite
calls ``mainly theoretic in nature and thus very complicated to
implement''~\citep{grafite2024} --- and measured its cost empirically.
The measurements show that adaptive binary/linear search on packed
suffixes outperforms WPS at LSM bucket sizes on both axes ---
$59$~bits per key of auxiliary metadata removed and per-query latency
reduced by $8$--$160\times$ relative to WPS.
Once WPS is gone, Goswami's two metadata vectors $D_1, D_2$ collapse
into a single Elias--Fano-encoded vector $D$ --- a further $24\%$
metadata saving and the same backend Grafite arrives at by
implementation-simplicity arguments (Chapter~\ref{chap:ere}).

On top of this backend we add a distribution-aware segmentation that
detects dense clusters and handles them exactly --- going below
Goswami's bound on real-world data (Search on Sorted Data
benchmark~\citep{kipf2019sosd}) while preserving the guarantee on
adversarial inputs.
